\chapter{Research-Grade categorate Enrichment}

\WeekHeader{5}{Research-Grade categorate Enrichment}{Use public database enrichment to convert chemical names into traceable research tables.}{Validated \texttt{categorate(detail = "research")} result and source coverage report.}

\section{Why This Matters}
The upgraded \texttt{categorate()} workflow gives students much more than a simple database lookup. It returns source-specific tables, normalized chemical traits, ontology-style terms, measurement summaries, hazard fields, occurrence evidence, KEGG roles, quality diagnostics, and data dictionaries. This is where the chemical list becomes a research dataset.

\section{Key Output Families}
\begin{itemize}
  \item \textbf{Identity and properties}: \texttt{PubChemIdentity}, \texttt{PubChemProperties}, \texttt{PubChemIdentifiers}.
  \item \textbf{Source profiles}: \texttt{SafetyProfile}, \texttt{FEMAProfile}, \texttt{FDA\_SPL\_Profile}, \texttt{LOTUSProfile}, \texttt{MeSHProfile}, \texttt{LiteratureProfile}.
  \item \textbf{Discrete trait data}: \texttt{ChemicalTerms}, \texttt{ChemicalTraits}, \texttt{ChemicalTraitOntology}.
  \item \textbf{Analysis matrices}: \texttt{ChemicalTraitMatrix}, \texttt{ChemicalTraitOntologyMatrix}.
  \item \textbf{Research summaries}: \texttt{ChemicalTraitReport}, \texttt{ChemicalMeasurementSummary}, \texttt{SourceCoverage}, \texttt{DerivedGroups}.
  \item \textbf{Database quality}: \texttt{DataDictionary}, \texttt{TableQuality}, \texttt{SourceDiagnostics}, \texttt{ValidationSummary}.
\end{itemize}

\section{Guided Lab}
Run a small enrichment first. Use caching for routine work and set a gentle throttle when live services are queried.

\begin{rcode}
library(uafR)
data("library_data", package = "uafR")

compounds <- c("aspirin", "caffeine", "methyl salicylate")

result <- categorate(
  compounds = compounds,
  chemical_library = library_data,
  input_format = "wide",
  detail = "research",
  cache = TRUE,
  throttle = 0.2,
  assay_detail_limit = 0,
  trait_matrix_profile = "core",
  trait_matrix_min_confidence = "medium",
  trait_matrix_max_traits = 100
)

validateCategorateResult(result)$Summary
result$SourceCoverage
result$ChemicalTraitReport
result$ChemicalTraitMatrix
\end{rcode}

\section{Evidence Inspection}
Students should not trust a matrix until they inspect the supporting evidence.

\begin{rcode}
trait_key <- names(result$ChemicalTraitMatrix)[2]

evidence <- chemicalTraitEvidence(result)
evidence[evidence$MatrixKey == trait_key, ]
\end{rcode}

\section{Independent Extension}
Students run \texttt{categorate()} on their reviewed chemical list. If the list is long, they should begin with 5 to 10 representative chemicals, inspect output quality, then scale up.

\section{Deliverables}
\begin{tabularx}{\textwidth}{p{0.25\textwidth}YY}
\DeliverableTableHeader
Enrichment object & Saved result contains validation tables, source coverage, and trait matrices. & RDS\\
Coverage report & Summarizes which sources returned useful data by chemical. & CSV or table\\
Evidence audit & Traces five matrix traits back to source evidence. & Research log\\
\end{tabularx}

\section{Quality Checklist}
\begin{checklistbox}{Before Week 6}
\begin{itemize}
  \item \texttt{ValidationSummary} has been inspected.
  \item Missing source data are documented rather than silently ignored.
  \item At least five high-value traits have been traced to evidence rows.
  \item The student can identify which output tables are raw-ish source evidence and which are normalized analysis tables.
\end{itemize}
\end{checklistbox}
